\documentclass[11pt]{article}
\usepackage[margin=25mm]{geometry}
\usepackage{amsmath,amssymb}
\usepackage{amsthm}
\newtheorem{proposition}{Proposition}[section]
\DeclareMathOperator{\Revenue}{Rev}
\DeclareMathOperator{\Cost}{Cost}
\begin{document}
\title{A guided mathematics example}
\author{}
\date{}
\maketitle
\section{Define the notation}
Let $q\in\mathbb{R}$ denote a nonnegative output quantity.
The unit price is $p\ge0$. The cost parameters satisfy $a>0$ and
$b>0$. This is a hypothetical model for practising notation.

\section{Group the related equations}
Revenue and cost are given by
\begin{subequations}\label{eq:model}
\begin{align}
\Revenue(q)&=pq,\label{eq:revenue}\\
\Cost(q)&=aq+bq^2.\label{eq:cost}
\end{align}
\end{subequations}
System~\eqref{eq:model} contains the revenue equation
\eqref{eq:revenue} and the cost equation~\eqref{eq:cost}.

Define profit $\Pi(q)$ as revenue minus cost. One derivation has
one number:
\begin{equation}\label{eq:profit}
\begin{aligned}
\Pi(q)&=\Revenue(q)-\Cost(q)\\
&=pq-(aq+bq^2)\\
&=(p-a)q-bq^2.
\end{aligned}
\end{equation}
The equals signs in \eqref{eq:profit} line up at the alignment markers.

\section{Write a boundary condition}
For an interior optimum the derivative satisfies
\[
\Pi'(q)=p-a-2bq=0.
\]
Since $\Pi''(q)=-2b<0$ and $q\ge0$, the maximizing quantity is
\[
q^*=\begin{cases}
\dfrac{p-a}{2b}&\text{if }p>a,\\[3pt]
0&\text{if }p\le a.
\end{cases}
\]
The two conditions cover the boundary $p=a$ exactly once.

\begin{proposition}\label{prop:nonnegative-cost}
For every $q\ge0$, the cost in \eqref{eq:cost} is nonnegative.
\end{proposition}
\begin{proof}
The assumptions imply $aq\ge0$ and $bq^2\ge0$.
Their sum is nonnegative.
\end{proof}
Proposition~\ref{prop:nonnegative-cost} follows from the stated signs,
rather than from any typesetting command.
\end{document}
